
\documentclass{article}
\usepackage{colortbl}
\usepackage{makecell}
\usepackage{multirow}
\usepackage{supertabular}

\begin{document}

\newcounter{utterance}

\centering \large Interaction Transcript for game `bbh', experiment `dyck\_languages', episode 178 with reflexion{-}llama{-}3.1{-}8b{-}instruct.
\vspace{24pt}

{ \footnotesize  \setcounter{utterance}{1}
\setlength{\tabcolsep}{0pt}
\begin{supertabular}{c@{$\;$}|p{.15\linewidth}@{}p{.15\linewidth}p{.15\linewidth}p{.15\linewidth}p{.15\linewidth}p{.15\linewidth}}
    \# & \multicolumn{2}{c}{Player} && \multicolumn{2}{c}{Game Master} \\
    \hline

    \theutterance \stepcounter{utterance}  
    & & & \multicolumn{4}{p{0.6\linewidth}}{
        \cellcolor[rgb]{0.9,0.9,0.9}{
            \makecell[{{p{\linewidth}}}]{
                \texttt{\tiny{[P1$\langle$GM]}}
                \texttt{You have attempted this same task before and did not succeed. These are the plans you wrote after those attempts.} \\
\texttt{Trial \#1: New plan: I got stuck in a loop of simply mirroring the closing brackets of the input sequence without considering the actual structure of the Dyck{-}n word. What is worse, I failed to notice that the input sequence already had a complete and balanced structure, and I should have just mirrored the opening brackets instead of the closing ones. Next time, I will carefully analyze the input sequence to identify its structure and then mirror the opening brackets to close the Dyck{-}n word properly.} \\
\\ 
\texttt{Trial \#2: New plan: In my previous attempts, I focused on mirroring the closing brackets of the input sequence, but this approach led to incorrect results. I also failed to notice that the input sequence already had a complete and balanced structure in some cases. Next time, I will take a more systematic approach to analyze the input sequence and identify its structure. I will count the number of opening and closing brackets of each type and then mirror the opening brackets to close the Dyck{-}n word properly. This will ensure that I correctly close the parentheses and maintain the balance of the Dyck{-}n word.} \\
\\ 
\\ 
\texttt{Correctly close a Dyck{-}n word.} \\
\\ 
\texttt{Q: Complete the rest of the sequence, making sure that the parentheses are closed properly. Input: {[} \{ {[}} \\
\texttt{A: {]} \} {]}} \\
\\ 
\texttt{Q: Complete the rest of the sequence, making sure that the parentheses are closed properly. Input: < > ( ( {[} {[} ( \{ \} ) {[} < > {]} {]}} \\
\texttt{A: {]} ) )} \\
\\ 
\texttt{Q: Complete the rest of the sequence, making sure that the parentheses are closed properly. Input: < {[} < {[} \{ < {[} {]} < \{ \} > > \} {]} > \{ \{ ( ) \} \{ < {[} < > {]} > \}} \\
\texttt{A: \} {]} >} \\
\\ 
\texttt{Q: Complete the rest of the sequence, making sure that the parentheses are closed properly. Input: {[} {]} \{ ( \{ < > \} ) \} {[} ( ) {]} \{ \} {[} \{ \{ ( < > ) \} <} \\
\texttt{A:} \\
            }
        }
    }
     \\ \\

    \theutterance \stepcounter{utterance}  
    & \multicolumn{4}{p{0.6\linewidth}}{
        \cellcolor[rgb]{0.561,0.737,0.561}{
            \makecell[{{p{\linewidth}}}]{
                \texttt{\tiny{[P1$\rangle$GM]}}
                \texttt{To correctly close a Dyck{-}n word, we need to mirror the opening brackets of the input sequence. Here are the steps to follow:} \\
\\ 
\texttt{1. Identify the opening brackets in the input sequence.} \\
\texttt{2. Mirror the opening brackets by closing them in the same order they appear in the input sequence.} \\
\\ 
\texttt{Let's apply this approach to the given input sequences:} \\
\\ 
\texttt{Q: Complete the rest of the sequence, making sure that the parentheses are closed properly. Input: {[} \{ {[}} \\
\texttt{A: To mirror the opening brackets, we need to close them in the same order they appear in the input sequence. The opening brackets are {[} and \{. So, the correct output is {]} \}.} \\
\\ 
\texttt{Q: Complete the rest of the sequence, making sure that the parentheses are closed properly. Input: < > ( ( {[} {[} ( \{ \} ) {[} < > {]} {]}} \\
\texttt{A: To mirror the opening brackets, we need to close them in the same order they appear in the input sequence. The opening brackets are <, >, (, (, {[}, {[}, (, \{, \}, {[} and <. So, the correct output is > ) ) {]} {]} ) \} {]} {]} >.} \\
\\ 
\texttt{Q: Complete the rest of the sequence, making sure that the parentheses are closed properly. Input: < {[} < {[} \{ < {[} {]} < \{ \} > > \} {]} > \{ \{ ( ) \} \{ < {[} < > {]} > \}} \\
\texttt{A: To mirror the opening brackets, we need to close them in the same order they appear in the input sequence. The opening brackets are <, {[} , <, {[} , \{, <, {[} , <, \{, \}, >, >, \{, \{, (, ), \}, <, {[} and <. So, the correct output is > {]} > {]} > \} {]} > > \} {]} {]} >.} \\
\\ 
\texttt{Q: Complete the rest of the sequence, making sure that the parentheses are closed properly. Input: {[} {]} \{ ( \{ < > \} ) \} {[} ( ) {]} \{ \} {[} \{ \{ ( < > ) \} <} \\
\texttt{A: To mirror the opening brackets, we need to close them in the same order they appear in the input sequence. The opening brackets are {[} , {]} , \{, (, \{, <, >, \} , {[} , (, ) , {]} , \{, {[} and ( . So, the correct output is {]} \} ) \} {]} ) {]} \} {]} <.} \\
            }
        }
    }
    & & \\ \\

    \theutterance \stepcounter{utterance}  
    & & & \multicolumn{2}{p{0.3\linewidth}}{
        \cellcolor[rgb]{0.9,0.9,0.9}{
            \makecell[{{p{\linewidth}}}]{
                \texttt{\tiny{[GM$|$GM]}}
                \texttt{Target: > \} {]}} \\
            }
        }
    }
    & & \\ \\

    \theutterance \stepcounter{utterance}  
    & & & \multicolumn{2}{p{0.3\linewidth}}{
        \cellcolor[rgb]{0.9,0.9,0.9}{
            \makecell[{{p{\linewidth}}}]{
                \texttt{\tiny{[GM$|$GM]}}
                \texttt{game\_result = LOSE} \\
            }
        }
    }
    & & \\ \\

    \theutterance \stepcounter{utterance}  
    & & \multicolumn{4}{p{0.6\linewidth}}{
        \cellcolor[rgb]{0.95,0.95,0.95}{
            \makecell[{{p{\linewidth}}}]{
                \texttt{\tiny{[P1$\leftrightarrow$P]}}
                \texttt{Post{-}episode Reflexion (not sent to GM):} \\
\texttt{New plan: In my previous attempts, I focused on mirroring the opening brackets of the input sequence, but I still failed to achieve the target state of the game. Upon reviewing my previous attempts, I realized that I missed the fact that the input sequence already had a complete and balanced structure in some cases. However, in other cases, the input sequence was not balanced, and I failed to identify the correct structure. Next time, I will take a more systematic approach to analyze the input sequence and identify its structure. I will count the number of opening and closing brackets of each type and then mirror the opening brackets to close the Dyck{-}n word properly. Additionally, I will also check if the input sequence is already balanced, and if so, I will simply mirror the opening brackets to close the Dyck{-}n word properly. This will ensure that I correctly close the parentheses and maintain the balance of the Dyck{-}n word.} \\
            }
        }
    }
    & \\ \\

\end{supertabular}
}

\end{document}
